% संपूर्ण मराठी ग्रंथातील प्रकरण 3: फलने
% हा संपादनयोग्य स्रोतखंड आहे. संकलनासाठी संपूर्ण स्रोतसंग्रहातील
% अक्षररूपे, पूर्वभाग, चिन्हव्याख्या आणि आकृती-स्रोत आवश्यक आहेत.
% मूळ: Open Logic Project; मजकूर आणि रूपांतर CC BY 4.0.
% यंत्रानुवाद, दुरुस्ती आणि तपासणी: OpenAI Codex —
% GPT-5.6 Sol आणि GPT-6 Sol, दोन्ही Ultra effort.
% दुरुस्ती-पुनरावलोकन: GPT-6.1 Sol, Ultra effort.
\chapter{फलने}\label{sfr:fun::chap}





\section{मूलभूत संकल्पना}\label{sfr:fun:bas:sec}

\begin{explain}
दिलेल्या संचातील प्रत्येक घटक दुसऱ्या एखाद्या दिलेल्या संचातील
एका ठरावीक घटक शी जोडणारी संगती म्हणजे \emph{फलन}.
उदाहरणार्थ, $1$ मिळवण्याच्या क्रियेने एक फलन ठरते: प्रत्येक संख्या~$n$
एका एकमेव संख्या~$n+1$ शी जोडली जाते.

अधिक व्यापकपणे पाहता, फलनांना आदान म्हणून जोड्या, त्रिके इत्यादी देता
येतात आणि त्यांच्याकडून कोणत्या तरी प्रकारचे प्रदान मिळते. प्राथमिक
अंकगणितातून अनेक फलने आपल्याला परिचित आहेत. उदाहरणार्थ, बेरीज आणि
गुणाकार ही फलने आहेत. ती दोन संख्या आदान म्हणून घेतात आणि तिसरी संख्या देतात.

या गणिती, अमूर्त अर्थाने फलन म्हणजे एक \emph{काळी पेटी}: कोणत्या आदानाशी
कोणते प्रदान जोडलेले आहे एवढेच महत्त्वाचे असते; प्रदानाची गणना करण्याची
पद्धत महत्त्वाची नसते.
\end{explain}

\begin{defn}[फलन]
\emph{फलन} $f \colon A \to B$ ही~$A$ मधील प्रत्येक घटक ची~$B$
मधील एका घटक शी लावलेली संगती असते.

$A$ ला~$f$ चा \emph{प्रांत} आणि $B$ ला~$f$ चा \emph{सहप्रांत} म्हणतो.
~$A$ मधील घटक ना~$f$ ची आदाने किंवा \emph{फलसाधक} म्हणतात.
फलसाधक~$x$ शी~$f$ ने जोडलेल्या~$B$ मधील घटक ला फलसाधक~$x$
साठीचे \emph{$f$ चे मूल्य} म्हणतात आणि ते~$f(x)$ असे लिहितात.

~$f$ ची \emph{व्याप्ती} $\ran{f}$ म्हणजे सहप्रांताचा असा उपसंच की
त्याचे घटक कोणत्या तरी फलसाधकासाठीची~$f$ ची मूल्ये असतात;
$\ran{f} = \Setabs{f(x)}{x \in A}$.
\end{defn}

फलनांचा विचार करण्यासाठी \ref{sfr:fun:bas:fig:function} मधील आकृती उपयोगी पडू शकते.
डावीकडील लंबवर्तुळ फलनाचा \emph{प्रांत} दर्शवते; उजवीकडील लंबवर्तुळ
त्याचा \emph{सहप्रांत} दर्शवते; आणि बाण प्रांतातील \emph{फलसाधकापासून}
सहप्रांतातील त्याच्या संगत \emph{मूल्याकडे} जातो.

\begin{figure}
  \olasset{assets/diagrams/function.tikz}
  \caption{फलन एका संचातील प्रत्येक घटक ची दुसऱ्या संचातील
    घटक शी संगती लावते. बाण प्रांतातील फलसाधकापासून
    सहप्रांतातील त्याच्या संगत मूल्याकडे जातो.}
  \label{sfr:fun:bas:fig:function}
\end{figure}

\begin{ex}
गुणाकार नैसर्गिक संख्यांच्या जोड्या आदान म्हणून घेतो आणि त्यांना प्रदान
म्हणून नैसर्गिक संख्या जोडतो. म्हणून तो $\Nat \times \Nat$ या प्रांताकडून
$\Nat$ या सहप्रांताकडे जातो. त्याची व्याप्तीही $\Nat$ च असते, कारण
प्रत्येक $n \in \Nat$ हा $n \times 1$ असतो.
\end{ex}

\begin{ex}
गुणाकार हे फलन आहे, कारण तो प्रत्येक आदानाशी---नैसर्गिक संख्यांच्या
प्रत्येक जोडीशी---एकच प्रदान जोडतो: $\times \colon \Nat^2 \to \Nat$.
याउलट, नैसर्गिक संख्या~$n$ शी $x^2 = n$ असणारी वास्तव संख्या~$x$
जोडण्याची संगती फलन होत नाही, कारण प्रत्येक धन पूर्णांक~$n$ ची दोन
वर्गमुळे असतात: $\sqrt{n}$ आणि~$-\sqrt{n}$. फक्त धन वर्गमूळ देऊन
आपण तिचे फलन करू शकतो: फक्त ऋणेतर (प्रधान) वर्गमूळ परत द्यायचे,
$\sqrt{\phantom{X}} \colon \Nat \to \Real$.
\end{ex}
\begin{quote}\small\textbf{स्रोतदुरुस्ती.} OLFUN-002 — गोठवलेल्या इंग्रजी स्रोतामध्ये सर्व नैसर्गिक संख्यांसाठी positive square root म्हटले आहे. शून्याचे प्रधान वर्गमूळ शून्य असल्यामुळे येथे ऋणेतर (प्रधान) असे दुरुस्त केले आहे; धन पूर्णांकांना दोन वास्तव वर्गमुळे असतात हे आधीचे विधान योग्य आहे.\end{quote}

\begin{ex}
वर्गातील प्रत्येक विद्यार्थ्याशी त्याची अंतिम श्रेणी जोडणारा संबंध हे
फलन आहे---एकाच वर्गात कोणत्याही विद्यार्थ्याला दोन वेगवेगळ्या अंतिम
श्रेणी मिळू शकत नाहीत. वर्गातील प्रत्येक विद्यार्थ्याशी त्याचे पालक
जोडणारा संबंध फलन नाही: विद्यार्थ्यांचे पालक शून्य, दोन किंवा अधिक असू शकतात.
\end{ex}

\begin{explain}
प्रत्येक शक्य फलसाधकासाठी फलनाचे मूल्य काय असते हे नेमक्या पद्धतीने
सांगून फलनाची व्याख्या करता येते. सूत्र देणे, मूल्याची गणना करण्याची
पद्धत वर्णन करणे किंवा प्रत्येक फलसाधकासाठीची मूल्ये यादीत देणे या
त्याच्या वेगवेगळ्या पद्धती आहेत. फलनाची व्याख्या कोणत्याही पद्धतीने
केली, तरी प्रत्येक फलसाधकासाठी एक आणि केवळ एकच मूल्य दिलेले असेल
याची खात्री केली पाहिजे.
\end{explain}


\begin{ex}
$f \colon \Nat \to \Nat$ ची व्याख्या $f(x) = x+1$ अशी करू.
या व्याख्येत $f$ हे नैसर्गिक संख्या आदान म्हणून घेऊन नैसर्गिक संख्या
प्रदान करणारे फलन आहे असे सांगितले आहे. नैसर्गिक संख्या~$x$ दिली,
की $f$ तिची उत्तरवर्ती संख्या~$x+1$ देईल असे ती सांगते.
या उदाहरणात सहप्रांत $\Nat$ हा~$f$ ची व्याप्ती नाही, कारण नैसर्गिक
संख्या~$0$ ही कोणत्याही नैसर्गिक संख्येची उत्तरवर्ती संख्या नाही.
~$f$ ची व्याप्ती सर्व धन पूर्णांकांचा संच $\Int^{+}$ आहे.
\end{ex}

\begin{ex}\label{sfr:fun:bas:examplefunext}
$g \colon \Nat \to \Nat$ ची व्याख्या $g(x) = x+2-1$ अशी करू.
यावरून $g$ हे नैसर्गिक संख्या आदान म्हणून घेऊन नैसर्गिक संख्या प्रदान
करणारे फलन आहे हे कळते. नैसर्गिक संख्या x दिली, की $g$ हा~$x$ च्या
उत्तरवर्ती संख्येच्या उत्तरवर्ती संख्येची पूर्ववर्ती संख्या,
म्हणजेच $x+1$, देईल.
\end{ex}
\begin{quote}\small\textbf{स्रोतदुरुस्ती.} OLFUN-003 — गोठवलेल्या इंग्रजी स्रोताच्या गद्यामध्ये याच उदाहरणात आदानासाठी आधी $n$ आणि लगेच x वापरले आहेत. सूत्र न बदलता मराठी गद्याने x हे एकच नाव सातत्याने वापरले आहे.\end{quote}

\begin{explain}
आपण नुकतीच वेगवेगळ्या \emph{व्याख्या} असलेली $f$ आणि $g$ ही दोन
फलने पाहिली. पण ती \emph{एकच फलन} आहेत. कारण कोणत्याही नैसर्गिक
संख्या~$n$ साठी $f(n) = n+1 = n+2-1 = g(n)$ असते. दुसऱ्या शब्दांत,
$f$ आणि~$g$ यांच्या व्याख्या वेगवेगळ्या समीकरणांनी एकच संगती सांगतात.
म्हणून आपण अप्रत्यक्षपणे फलनांसाठीच्या विस्तारात्मकतेच्या तत्त्वावर
अवलंबून आहोत:
\[  \text{जर }\forall x\, f(x) = g(x)\text{, तर }f = g
\]
मात्र $f$ आणि~$g$ यांचा प्रांत आणि सहप्रांत एकच असला पाहिजे.
\end{explain}

\begin{ex}
प्रकरणे पाडूनही फलनाची व्याख्या करता येते. उदाहरणार्थ,
$h \colon \Nat \to \Nat$ ची व्याख्या अशी करता येईल:
\[h(x) =
\begin{cases}
  \frac{x}{2} & \text{जर $x$ सम असेल} \\
  \frac{x+1}{2} & \text{जर $x$ विषम असेल.}
\end{cases}
\]
प्रत्येक नैसर्गिक संख्या सम किंवा विषम असते, म्हणून या फलनाचे प्रदान
नेहमी नैसर्गिक संख्या असेल. प्रकरणे पाडून फलनाची व्याख्या करताना
प्रत्येक शक्य आदान नेमक्या एका प्रकरणात बसले पाहिजे, हे लक्षात ठेवा.
काही वेळा सर्व शक्यता त्या प्रकरणांत येतात आणि कोणतीही शक्यता दोन
प्रकरणांत येत नाही याची सिद्धता करावी लागेल.
\end{ex}








\section{फलनांचे प्रकार}\label{sfr:fun:kin:sec}

\begin{explain}
आपल्याला वारंवार भेटणाऱ्या फलनांच्या काही प्रकारांचे वर्गीकरण करून
घेणे उपयोगी ठरेल.

सुरुवातीला, सहप्रांतातील प्रत्येक सदस्य फलनाचे मूल्य असतो असा गुणधर्म
असलेल्या फलनांचा विचार करू. अशा फलनांना आच्छादक म्हणतात.
त्यांची आकृती \ref{sfr:fun:kin:fig:surjective} प्रमाणे काढता येते.

\begin{figure}
  \olasset{assets/diagrams/surjective.tikz}
  \caption{आच्छादक फलनात सहप्रांतातील प्रत्येक घटक
    हे फलनाचे मूल्य असते.}
  \label{sfr:fun:kin:fig:surjective}
\end{figure}
\end{explain}

\begin{defn}[आच्छादक फलन]
$f \colon A \rightarrow B$ हे फलन \emph{आच्छादक} असते तेव्हा
आणि केवळ तेव्हाच, जेव्हा $B$ हा~$f$ ची व्याप्तीही असतो. म्हणजे,
प्रत्येक $y \in B$ साठी~$f(x) = y$ असणारा किमान एक $x \in A$
असतो. चिन्हांत:
\[  (\forall y \in B)(\exists x \in A)f(x) = y.
\]
अशा फलनाला $A$ पासून $B$ पर्यंतचे आच्छादन म्हणतो.
\end{defn}

\begin{explain}
$f$ हे आच्छादन आहे असे दाखवायचे असल्यास, $f$ च्या सहप्रांतातील
प्रत्येक वस्तू कोणत्या तरी आदान $x$ साठी $f(x)$ चे मूल्य आहे असे दाखवावे लागते.

कोणत्याही फलनापासून एक आच्छादन \emph{निर्माण होते}, हे लक्षात घ्या.
कारण $f \colon A \to B$ हे फलन दिले असल्यास, $f' \colon A \to \ran{f}$
ची व्याख्या $f'(x) = f(x)$ अशी करू. $\ran{f}$ ची \emph{व्याख्याच}
$\Setabs{f(x) \in B}{x \in A}$ अशी असल्यामुळे हे फलन $f'$ निश्चितच
आच्छादन असते.
\end{explain}

\begin{explain}
कोणतेही फलन प्रत्येक शक्य आदानाशी एकमेव प्रदान जोडते. पण वेगवेगळ्या
आदानांशी कधीच एकच प्रदान न जोडणारी फलनेही असतात. अशा फलनांना
एकास-एक म्हणतात. त्यांची आकृती \ref{sfr:fun:kin:fig:injective} प्रमाणे काढता येते.
\begin{figure}
  \olasset{assets/diagrams/injective.tikz}
  \caption{एकास-एक फलन कधीच दोन वेगवेगळ्या फलसाधकांशी
    एकच मूल्य जोडत नाही.}
  \label{sfr:fun:kin:fig:injective}
\end{figure}
\end{explain}

\begin{defn}[एकास-एक फलन]
$f \colon A \rightarrow B$ हे फलन \emph{एकास-एक} असते तेव्हा
आणि केवळ तेव्हाच, जेव्हा प्रत्येक $y \in B$ साठी~$f(x) = y$ असणारा
जास्तीत जास्त एक $x \in A$ असतो. अशा फलनाला $A$ पासून~$B$ पर्यंतचे
एकास-एक फलन म्हणतो.
\end{defn}

\begin{explain}
$f$ हे एकास-एक फलन आहे असे दाखवायचे असल्यास, $f$ च्या प्रांतातील
कोणत्याही घटक $x$ आणि $y$ साठी $f(x)=f(y)$ असल्यास $x=y$
असते असे दाखवावे लागते.
\end{explain}

\begin{ex}
$f(x) = 1$ ने दिलेले स्थिर फलन $f\colon \Nat \to \Nat$ हे
एकास-एक ही नाही आणि आच्छादक ही नाही.

$f(x) = x$ ने दिलेले एकरूपता फलन $f\colon \Nat \to \Nat$ हे
एकास-एक आणि आच्छादक दोन्ही आहे.

$f(x) = x+1$ ने दिलेले उत्तरवर्ती फलन $f \colon \Nat \to \Nat$
हे एकास-एक आहे, पण आच्छादक नाही.

पुढीलप्रमाणे व्याख्या केलेले फलन $f \colon \Nat \to \Nat$:
\[  f(x) =
  \begin{cases}
    \frac{x}{2} & \text{जर $x$ सम असेल} \\
    \frac{x+1}{2} & \text{जर $x$ विषम असेल.}
  \end{cases}
\]
हे आच्छादक आहे, पण एकास-एक नाही.
\end{ex}

\begin{explain}
अनेकदा एकास-एक आणि आच्छादक हे दोन्ही गुणधर्म असलेल्या
फलनांचा विचार करायचा असतो. अशा फलनांना एकास-एक व आच्छादक म्हणतो.
ती \ref{sfr:fun:kin:fig:bijective} मधील फलनासारखी दिसतात. एकास-एक आच्छादने ना
कधी कधी \emph{एकास-एक संगती} असेही म्हणतात, कारण ती सहप्रांतातील
घटकांची प्रांतातील घटकांशी एकमेव जोडी लावतात.
\begin{figure}
  \olasset{assets/diagrams/bijective.tikz}
  \caption{एकास-एक व आच्छादक फलन सहप्रांतातील घटकांची प्रांतातील
    घटकांशी एकमेव जोडी लावते.}
  \label{sfr:fun:kin:fig:bijective}
\end{figure}
\end{explain}

\begin{defn}[एकास-एक आच्छादन]
$f \colon A \to B$ हे फलन \emph{एकास-एक व आच्छादक} असते तेव्हा आणि
केवळ तेव्हाच, जेव्हा ते आच्छादक आणि एकास-एक दोन्ही असते.
अशा फलनाला $A$ पासून~$B$ पर्यंतचे (किंवा $A$ आणि~$B$ यांमधील)
एकास-एक आच्छादन म्हणतो.
\end{defn}










\section{संबंध म्हणून फलने}\label{sfr:fun:rel:sec}

\begin{explain}
~$A$ मधील घटक ची~$B$ मधील घटक शी संगती लावणारे
फलन $A$ आणि~$B$ यांच्यामधील एक संबंध ठरवते: $f(x) = y$ असेल
तेव्हा आणि केवळ तेव्हाच $x$ आणि $y$ यांच्यात हा संबंध असतो.
खरे तर, उदाहरणार्थ गणिताच्या पायाभूत घटकांची संख्या कमी करायची असेल,
तर आपण फलन~$f$ आणि हा संबंध, म्हणजे जोड्यांचा एक संच, \emph{एकरूप}
मानू शकतो. मग असा प्रश्न उभा राहतो: कोणते संबंध अशा रीतीने फलने ठरवतात?
\end{explain}

\begin{defn}[फलनाचा आलेख] $f\colon A \to B$ हे फलन असू द्या.
~$f$ चा \emph{आलेख} म्हणजे पुढील व्याख्येने मिळणारा संबंध
$R_f \subseteq A \times B$:
\[R_f = \Setabs{\tuple{x,y}}{f(x) = y}.
\]
\end{defn}

\begin{explain}
विस्तारात्मकतेमुळे फलनाचा आलेख एकमेव ठरतो. शिवाय, संचांसाठीच्या
विस्तारात्मकतेमुळे फलनांसाठीचे अप्रत्यक्ष विस्तारात्मकता तत्त्व लगेच
समर्थित होते: $f$ आणि~$g$ यांचा प्रांत व सहप्रांत एकच असेल आणि
सर्व मूल्यांवर ती जुळत असतील, तर ती एकच फलन असतात.

तसेच, एखादा संबंध ``फलनरूप'' असेल, तर तो एका फलनाचा आलेख असतो.
\end{explain}

\begin{prop}\label{sfr:fun:rel:prop:graph-function}
$R \subseteq A \times B$ हा संबंध पुढील अटी पूर्ण करतो असे समजा:
\begin{enumerate}
\item $Rxy$ आणि $Rxz$ असल्यास $y = z$ असते; आणि
\item प्रत्येक $x \in A$ साठी $\tuple{x,
y} \in R$ असणारा काही $y \in B$ असतो.
\end{enumerate}
मग $f(x) = y$ तेव्हा आणि केवळ तेव्हाच $Rxy$, अशा व्याख्येने मिळणाऱ्या
फलन $f\colon A \to B$ चा आलेख $R$ असतो.
\end{prop}

\begin{proof}
$Rxy$ असणारा एक $y$ आहे असे समजा. $Rxz$ असणारा दुसरा $z \neq y$
असता, तर~$R$ वरील अट मोडली गेली असती. म्हणून $Rxy$ असणारा एक $y$
असेल, तर हा $y$ एकमेव असतो. त्यामुळे $f$ सुस्पष्टपणे परिभाषित आहे.
उघडच, $R_f = R$.
\end{proof}

\begin{explain}
प्रत्येक फलन $f\colon A \to B$ ला एक आलेख असतो, म्हणजे $f(x) = y$
ने परिभाषित होणारा A आणि B यांच्यातील, $A \times B$ मध्ये
समाविष्ट असलेला एक संबंध असतो. उलट,
\ref{sfr:fun:rel:prop:graph-function} मधील गुणधर्म असलेला प्रत्येक संबंध~$R
\subseteq A \times B$ हा एका फलन~$f \colon A \to B$ चा आलेख असतो.
फलने आणि त्यांचे आलेख यांचा हा निकटचा संबंध असल्यामुळे फलन म्हणजे
त्याचा आलेखच असे आपण मानू शकतो. दुसऱ्या शब्दांत, फलने विशिष्ट
संबंधांशी, म्हणजे विशिष्ट क्रमित घटकसमूहांच्या संचांशी, एकरूप मानता
येतात. पण या ``एकरूप मानण्याचा'' आशय
\ref{sfr:rel:ref:sec} मधल्यासारखाच आहे, हे लक्षात घ्या: तो
फलनांच्या सत्तामीमांसेविषयीचा दावा नाही, तर फलनांना विशिष्ट संच
\emph{मानणे} सोयीचे आहे हे निरीक्षण आहे. ही सोय होण्याचे एक कारण असे की,
संबंधांवर जशा क्रिया केल्या तशाच क्रिया फलनांवर करण्याचा आता आपण
विचार करू शकतो (\ref{sfr:rel:ops:sec} पाहा). विशेषतः:
\end{explain}
\begin{quote}\small\textbf{स्रोतदुरुस्ती.} OLFUN-004 — इंग्रजी स्रोत येथे graph ला relation on A times B म्हणतो; पण त्याच ग्रंथाच्या व्याख्येनुसार त्याचा शब्दशः अर्थ (A times B) च्या वर्गाचा उपसंच होईल. मराठी मजकुरात अचूक अर्थ A आणि B यांच्यातील, A times B मध्ये समाविष्ट संबंध असा दिला आहे.\end{quote}

\begin{defn}\label{sfr:fun:rel:defn:funimage}
$f \colon A \to B$ हे फलन असू द्या आणि $C\subseteq A$ असू द्या.

~$f$ चे~$C$ पर्यंतचे \emph{मर्यादन} म्हणजे सर्व $x \in C$ साठी
$(\funrestrictionto{f}{C})(x) = f(x)$ या व्याख्येने मिळणारे
फलन~$\funrestrictionto{f}{C}\colon C \to B$. दुसऱ्या शब्दांत,
$\funrestrictionto{f}{C} = \Setabs{\tuple{x, y} \in R_f}{x \in C}$.

~$f$ चे~$C$ वरील \emph{उपयोजन} म्हणजे $\funimage{f}{C} =
\Setabs{f(x)}{x \in C}$. यालाच~$f$ अंतर्गत~$C$ ची \emph{प्रतिमा}
असेही म्हणतो.
\end{defn}

\begin{explain}
या व्याख्यांवरून कोणत्याही फलन~$f$ साठी $\ran{f} =
\funimage{f}{\dom{f}}$ असते.
\ref{sfr:rel:ops:sec} मधील
व्याख्या आणि फलनांना संबंधांशी एकरूप मानण्याची आपली पद्धत लक्षात
घेतल्यास या संकल्पना अनुरूप आहेत. मात्र फलनाचे C पर्यंतचे मर्यादन
केवळ आदान निर्देशांक C मध्ये ठेवते; संबंधाचे C वरील मर्यादन
दोन्ही निर्देशांक C मध्ये ठेवते. त्यामुळे त्या अगदी एकच क्रिया नाहीत.
व्युत्क्रम आणि सापेक्ष गुणाकार या इतर दोन क्रियांना थोडे अधिक
स्पष्टीकरण लागते. ते आपण
\ref{sfr:fun:inv:sec} आणि \ref{sfr:fun:cmp:sec} मध्ये देऊ.
\end{explain}
\begin{quote}\small\textbf{स्रोतदुरुस्ती.} OLFUN-005 — फलनाच्या मर्यादनाचे स्पष्ट सूत्र योग्य आहे: आलेखात पहिला निर्देशांक C मध्ये असतो. आधीचे संबंधावरील मर्यादन दोन्ही निर्देशांक C मध्ये ठेवते. इंग्रजी स्रोताचा exactly असा दावा येथे अनुरूप पण भिन्न असा स्पष्ट केला आहे.\end{quote}








\section{फलनांचे व्युत्क्रम}\label{sfr:fun:inv:sec}

\begin{explain}
फलनांचा विचार आपण संगती म्हणून करतो. मग ही संगती ``उलटवता'' येते का,
हा प्रश्न स्वाभाविकपणे पडतो. उदाहरणार्थ, उत्तरवर्ती फलन $f(x) = x + 1$
उलटवता येते, या अर्थाने की $g(y) = y - 1$ हे फलन $f$ ने केलेली
क्रिया ``पूर्ववत करते''.

पण येथे काळजी घेतली पाहिजे. ~$g$ च्या व्याख्येने $\Int \to \Int$
असे फलन ठरत असले, तरी $\Nat \to \Nat$ असे \emph{फलन} ठरत नाही,
कारण $g(0) \notin \Nat$. म्हणून साध्या उदाहरणांतसुद्धा फलन उलटवता
येते का हे तितकेसे उघड नसते; ते प्रांत आणि सहप्रांतावर अवलंबून असू शकते.

फलनाच्या व्युत्क्रमाच्या संकल्पनेने हा विचार अधिक नेमका करता येतो.
\end{explain}

\begin{defn}
सर्व $x \in A$ आणि $y \in B$ साठी $f(g(y)) = y$ आणि $g(f(x)) = x$
असेल, तर $g \colon B \to A$ हे फलन $f \colon A \to B$ या फलनाचा
\emph{व्युत्क्रम} असते.
\end{defn}

$f$ ला व्युत्क्रम~$g$ असेल, तर आपण अनेकदा~$g$ ऐवजी $f^{-1}$ असे लिहितो.

\begin{explain}
आता कोणत्या फलनांना व्युत्क्रम असतात हे ठरवू. $f\colon A \to B$
च्या व्युत्क्रमासाठी एक संभाव्य फलन म्हणजे पुढीलप्रमाणे ``व्याख्या केलेले''
$g\colon B \to A$:
\[g(y) = \text{$f(x) = y$ असणारा ``तो'' $x$.}
\]
पण ``व्याख्या केलेले'' आणि ``तो'' यांभोवतीची अवतरणचिन्हे सुचवतात की
ही खरे तर व्याख्या नाही. किमान पूर्ण व्यापकतेने ती नेहमी लागू पडणार
नाही. कारण या व्याख्येने फलन ठरवायचे असेल, तर $f(x) = y$ असणारा
एक आणि केवळ एकच~$x$ असला पाहिजे---~$g$ चे प्रदान एकमेव ठरले पाहिजे.
शिवाय ते प्रत्येक $y \in B$ साठी ठरले पाहिजे. $x_1$ आणि $x_2 \in A$
हे $x_1 \neq x_2$ असूनही $f(x_1) = f(x_2)$ असतील, तर
$y = f(x_1) = f(x_2)$ साठी $g(y)$ एकमेव ठरणार नाही.
आणि $f(x) = y$ असणारा~$x$ मुळीच नसेल, तर $g(y)$ मुळीच ठरणार नाही.
दुसऱ्या शब्दांत, $g$ परिभाषित होण्यासाठी $f$ हे एकास-एक आणि
आच्छादक दोन्ही असले पाहिजे.

हे टप्प्याटप्प्याने पाहू. प्रश्नाचे दोन भाग करू: फलन~$f\colon A \to B$
दिले असता, $g(f(x)) = x$ असणारे फलन $g\colon B \to A$ कधी असते?
असे $g$ हे $f$ ने केलेली क्रिया ``पूर्ववत करते'', आणि त्याला~$f$ चा
\emph{डावा व्युत्क्रम} म्हणतात. दुसरे म्हणजे, $f(h(y)) = y$ असणारे
फलन $h\colon B \to A$ कधी असते? अशा $h$ ला~$f$ चा
\emph{उजवा व्युत्क्रम} म्हणतात---$f$ हे~$h$ ने केलेली क्रिया ``पूर्ववत करते''.
\end{explain}

\begin{prop}
A अरिक्त असेल आणि $f\colon A \to B$ हे एकास-एक असेल,
तर~$f$ चा एक
\emph{डावा व्युत्क्रम}~$g\colon B \to A$ असतो, आणि सर्व $x \in A$
साठी $g(f(x)) = x$ असते.
\end{prop}
\begin{quote}\small\textbf{स्रोतदुरुस्ती.} OLFUN-001 — गोठवलेल्या इंग्रजी प्रमेयात A अरिक्त असण्याची अट सुटली आहे. रिक्त A पासून अरिक्त B कडेचे रिक्त फलन एकास-एक असूनही B पासून A कडे फलन नसते. नेमकी सामान्य अट: एकास-एक फलनाला डावा व्युत्क्रम तेव्हा आणि केवळ तेव्हाच असतो, जेव्हा A अरिक्त असतो किंवा B रिक्त असतो. मराठी प्रमेयात साधी पुरेशी A अरिक्त ही अट जोडली आहे.\end{quote}

\begin{proof}
A अरिक्त आहे आणि $f\colon A \to B$ हे एकास-एक आहे असे समजा.
$y \in B$ विचारात घ्या.
$y \in \ran{f}$ असेल, तर $f(x) = y$ असणारा एक $x \in A$ असतो.
$f$ हे एकास-एक असल्यामुळे असा~$x \in A$ एकच असतो. म्हणून
$g(y) = x$ अशी व्याख्या करता येते; म्हणजे $g(y)$ हा $f(x) = y$
असणारा ``तो'' $x \in A$ असतो. $y \notin \ran{f}$ असेल, तर त्याची
संगती कोणत्याही~$a \in A$ शी लावता येते. म्हणून एक $a \in A$
निवडून $g\colon B \to A$ ची व्याख्या अशी करता येते:
\[g(y) = \begin{cases}
    x & \text{जर $f(x) = y$ असेल}\\
    a & \text{जर $y \notin \ran{f}$ असेल.}
\end{cases}
\]
हे सर्व $y \in B$ साठी परिभाषित आहे, कारण अशा प्रत्येक $y \in \ran{f}$
साठी $f(x) = y$ असणारा नेमका एक $x \in A$ असतो. व्याख्येनुसार,
$y = f(x)$ असल्यास $g(y) = x$, म्हणजे $g(f(x)) = x$.
\end{proof}

\begin{prob}
$f\colon A \to B$ ला डावा व्युत्क्रम~$g$ असल्यास $f$ हे
एकास-एक असते, हे दाखवा.
\end{prob}

\begin{prop}
    $f\colon A \to B$ हे आच्छादक असेल, तर~$f$ चा एक
    \emph{उजवा व्युत्क्रम}~$h\colon B \to A$ असतो, आणि सर्व~$y \in B$
    साठी $f(h(y)) = y$ असते.
\end{prop}

\begin{proof}
$f\colon A \to B$ हे आच्छादक आहे असे समजा. $y \in B$ विचारात घ्या.
~$f$ हे आच्छादक असल्यामुळे $f(x_y) = y$ असणारा एक $x_y \in A$
असतो. म्हणून $h(y) = x_y$ अशी व्याख्या करता येते: प्रत्येक $y \in B$
साठी $f(x) = y$ असणारा काही $x \in A$ आपण निवडतो. ~$f$ हे
आच्छादक असल्यामुळे निवडण्यासाठी नेहमी किमान एक घटक असतो.\footnote{
$f$ हे आच्छादक असल्यामुळे प्रत्येक~$y \in B$ साठी
$\Setabs{x}{f(x) = y}$ हा संच अरिक्त असतो. ~$h$ ची आपली व्याख्या
या प्रत्येक संचातून एकच $x$ निवडण्याची मागणी करते. हे नेहमी करता येते
ही गोष्ट खरे तर उघड नाही---या निवडी करता येतात हे एक स्वयंसिद्धक म्हणून
गृहीत धरले जाते. दुसऱ्या शब्दांत, ही प्रतिज्ञा तथाकथित निवडीचे स्वयंसिद्धक
गृहीत धरते. या प्रश्नाकडे आपण सध्या दुर्लक्ष करू.
मात्र अनेक विशिष्ट उदाहरणांत, जसे $A = \Nat$ असेल किंवा सांत असेल,
किंवा $f$ हे एकास-एक व आच्छादक असेल, तेव्हा निवडीचे स्वयंसिद्धक लागत नाही.
($f$ हे एकास-एक व आच्छादक असण्याच्या विशिष्ट उदाहरणात प्रत्येक $y \in B$
साठी $\Setabs{x}{f(x) = y}$ या संचात नेमका एक घटक असतो,
म्हणून निवड करण्याचा प्रश्नच उरत नाही.)}
व्याख्येनुसार, $x = h(y)$ असल्यास $f(x) = y$ असते; म्हणजे कोणत्याही
$y \in B$ साठी $f(h(y)) = y$.
\end{proof}

\begin{prob}
$f\colon A \to B$ ला उजवा व्युत्क्रम~$h$ असल्यास $f$ हे
आच्छादक असते, हे दाखवा.
\end{prob}

\begin{explain}
  मागील सिद्धतेतील कल्पना एकत्र केल्यावर प्रत्येक एकास-एक आच्छादन ला
  व्युत्क्रम असतो हे मिळते. म्हणजे~$f$ चा डावा आणि उजवा व्युत्क्रम
  असे दोन्ही असणारे एकच फलन असते.
\end{explain}

\begin{prop}\label{sfr:fun:inv:prop:bijection-inverse}
$f\colon A \to B$ हे एकास-एक व आच्छादक असेल, तर एक
फलन~$f^{-1}\colon B \to A$ असते, आणि सर्व $x \in A$ साठी
$f^{-1}(f(x)) = x$ तसेच सर्व $y \in B$ साठी $f(f^{-1}(y)) = y$ असते.
\end{prop}

\begin{proof}
सरावासाठी.
\end{proof}

\begin{prob}
\ref{sfr:fun:inv:prop:bijection-inverse} सिद्ध करा. ~$f^{-1}$ ची
व्याख्या करा, ते फलन आहे हे दाखवा आणि ते~$f$ चा व्युत्क्रम आहे हे
दाखवा. म्हणजे, सर्व $x \in A$ आणि $y \in B$ साठी
$f^{-1}(f(x)) = x$ आणि $f(f^{-1}(y)) = y$ हे दाखवा.
\end{prob}

\begin{explain}
व्युत्क्रम मिळवण्याची थोडी अधिक व्यापक पद्धत आहे. \ref{sfr:fun:kin:sec}
मध्ये आपण पाहिले की प्रत्येक फलन $f$ पासून सर्व $x \in A$ साठी
$f'(x) = f(x)$ असे ठेवून आच्छादन $f' \colon A \to \ran{f}$
मिळते. उघडच, ~$f$ हे एकास-एक असेल, तर~$f'$ हे एकास-एक व आच्छादक
असते. म्हणून \ref{sfr:fun:inv:prop:bijection-inverse} नुसार त्याला एकमेव
व्युत्क्रम असतो. संकेतलेखनात अगदी थोडी सवलत घेऊन आपण कधी कधी
$f'$ च्या व्युत्क्रमाला फक्त ``~$f$ चा व्युत्क्रम'' म्हणतो.
\end{explain}

\begin{prop}\label{sfr:fun:inv:prop:left-right}%
  $f\colon A \to B$ ला डावा व्युत्क्रम~$g$ आणि उजवा व्युत्क्रम~$h$
  असल्यास $h = g$ असते, हे दाखवा.
\end{prop}

\begin{proof}
  सरावासाठी.
\end{proof}

\begin{prob}
  \ref{sfr:fun:inv:prop:left-right} सिद्ध करा.
\end{prob}

\begin{prop}\label{sfr:fun:inv:prop:inverse-unique}
प्रत्येक फलन~$f$ ला जास्तीत जास्त एक व्युत्क्रम असतो.
\end{prop}

\begin{proof}
  $g$ आणि $h$ हे दोन्ही~$f$ चे व्युत्क्रम आहेत असे समजा. मग विशेषतः
  ~$g$ हा~$f$ चा डावा व्युत्क्रम आणि~$h$ हा उजवा व्युत्क्रम असतो.
  \ref{sfr:fun:inv:prop:left-right} नुसार $g = h$.
\end{proof}








\section{फलनांचे संयोजन}\label{sfr:fun:cmp:sec}

\begin{explain}
\ref{sfr:fun:inv:sec} मध्ये आपण पाहिले की
एकास-एक आच्छादन~$f$ चा व्युत्क्रम~$f^{-1}$ हाही एक फलन असतो.
फलनांवरील आणखी एक क्रिया म्हणजे संयोजन: $f$ आणि~$g$ या दोन
फलनांचे संयोजन करून, म्हणजे आधी $f$ आणि मग~$g$ उपयोजून,
एक नवे फलन परिभाषित करता येते. अर्थात, व्याप्ती आणि प्रांत यांचा मेळ
असेल तेव्हाच हे शक्य असते; म्हणजे~$f$ ची व्याप्ती ही~$g$ च्या प्रांताचा
उपसंच असली पाहिजे. फलनांवरील ही क्रिया
\ref{sfr:rel:ops:sec} मधील संबंधांवरील सापेक्ष गुणाकाराच्या क्रियेशी
समांतर आहे.

संयोजनाची कल्पना स्पष्ट करण्यासाठी आकृती उपयोगी पडू शकते.
\ref{sfr:fun:cmp:fig:composition} मध्ये $f \colon A \to B$ आणि $g \colon B \to C$
ही दोन फलने आणि त्यांचे संयोजन~$(\comp{f}{g})$ दाखवले आहे.
$(\comp{f}{g}) \colon A \to C$ हे फलन~$A$ मधील प्रत्येक घटक
शी~$C$ मधील घटक जोडते. $A$ मधील घटक शी~$C$
मधील कोणता घटक जोडायचा हे असे ठरवतो: आदान $x \in A$
दिले असता, प्रथम~$x$ वर फलन $f$ उपयोजा. त्यामुळे काही
$f(x) = y \in B$ हे प्रदान मिळेल. मग~$y$ वर फलन $g$ उपयोजा.
त्यामुळे काही $g(f(x)) = g(y) = z \in C$ हे प्रदान मिळेल.
\begin{figure}
  \olasset[2\olphotowidth]{assets/diagrams/composition.tikz}
  \caption{$f$ आणि~$g$ या दोन फलनांचे संयोजन $g \circ f$.}
  \label{sfr:fun:cmp:fig:composition}
\end{figure}
\end{explain}

\begin{defn}[संयोजन]
$f\colon A \to B$ आणि $g\colon B \to C$ ही फलने असू द्या.
$f$ चे~$g$ बरोबरचे \emph{संयोजन} म्हणजे $\comp{f}{g} \colon A \to C$,
जेथे $(\comp{f}{g})(x) = g(f(x))$.
\end{defn}

\begin{ex}
$f(x) = x + 1$ आणि $g(x) = 2x$ ही फलने विचारात घ्या.
$(\comp{f}{g})(x) = g(f(x))$ असल्यामुळे प्रत्येक आदान~$x$ साठी
प्रथम त्याची उत्तरवर्ती संख्या घेऊन नंतर तिला दोनने गुणले पाहिजे.
म्हणून त्यांचे संयोजन $(\comp{f}{g})(x) = 2(x+1)$ असे मिळते.
\end{ex}

\begin{prob}
$f \colon A \to B$ आणि $g \colon B \to C$ ही दोन्ही फलने
एकास-एक असल्यास $\comp{f}{g}\colon A \to C$ हे
एकास-एक असते, हे दाखवा.
\end{prob}

\begin{prob}
$f \colon A \to B$ आणि $g \colon B \to C$ ही दोन्ही फलने
आच्छादक असल्यास $\comp{f}{g}\colon A \to C$ हे
आच्छादक असते, हे दाखवा.
\end{prob}

\begin{prob}
$f \colon A \to B$ आणि $g \colon B \to C$ असे समजा.
$\comp{f}{g}$ चा आलेख $R_f \mid R_g$ आहे हे दाखवा.
\end{prob}










\section{अंशतः फलने}\label{sfr:fun:par:sec}

\begin{explain}
कधी कधी फलनाच्या व्याख्येत थोडी सवलत देणे उपयोगी ठरते: सर्व शक्य
आदानांसाठी फलनाचे प्रदान परिभाषित असलेच पाहिजे ही अट काढून टाकता
येते. अशा संगतींना \emph{अंशतः फलने} म्हणतात.
\end{explain}

\begin{defn}
\emph{अंशतः फलन} $f \colon A \pto B$ ही अशी संगती असते की ती~$A$
मधील प्रत्येक घटक शी~$B$ मधील जास्तीत जास्त एक घटक
जोडते. $f$ ने $x \in A$ शी~$B$ मधील एखादा घटक जोडला असेल, तर
$f(x)$ \emph{परिभाषित} आहे असे म्हणतो; अन्यथा ते \emph{अपरिभाषित}
असते. $f(x)$ परिभाषित असल्यास $f(x) \fdefined$ असे लिहितो;
अन्यथा $f(x) \fundefined$ असे लिहितो. अंशतः फलन~$f$ चा \emph{प्रांत}
म्हणजे ते परिभाषित असलेल्या~$A$ मधील घटकांचा उपसंच:
$\dom{f} = \Setabs{x \in A}{f(x) \fdefined}$.
\end{defn}

\begin{ex}
प्रत्येक फलन $f\colon A \to B$ हे अंशतः फलनही असते. ~$A$ वरील
सर्व ठिकाणी परिभाषित असलेल्या अंशतः फलनांना---म्हणजे ज्यांना आपण
आतापर्यंत फक्त फलन म्हणत आलो, त्यांना---\emph{सर्वत्र परिभाषित}
फलने असेही म्हणतात.
\end{ex}

\begin{ex}
$f(x) = 1/x$ ने दिलेले अंशतः फलन $f \colon \Real \pto \Real$
हे $x = 0$ साठी अपरिभाषित असते आणि इतर सर्व ठिकाणी परिभाषित असते.
\end{ex}

\begin{prob}
$f\colon A \pto B$ दिले असता, अंशतः फलन $g\colon B \pto A$ ची
व्याख्या अशी करा: कोणत्याही $y \in B$ साठी $f(x) = y$ असणारा एकमेव
$x \in A$ असेल, तर $g(y) = x$; अन्यथा $g(y) \fundefined$.
$f$ हे एकास-एक असल्यास सर्व $x \in \dom{f}$ साठी $g(f(x)) = x$
आणि सर्व $y \in \ran{f}$ साठी $f(g(y)) = y$ असते, हे दाखवा.
\end{prob}

\begin{defn}[अंशतः फलनाचा आलेख]
$f\colon A \pto B$ हे अंशतः फलन असू द्या. ~$f$ चा \emph{आलेख}
म्हणजे पुढील व्याख्येने मिळणारा संबंध $R_f \subseteq A \times B$:
\[R_f = \Setabs{\tuple{x,y}}{f(x) = y}.
\]
\end{defn}

\begin{prop}
$R \subseteq A \times B$ या संबंधाचा असा गुणधर्म आहे असे समजा की
$Rxy$ आणि $Rxy'$ असल्यास $y = y'$ असते. मग $R$ हा पुढीलप्रमाणे
व्याख्या केलेल्या अंशतः फलन $f\colon A \pto B$ चा आलेख असतो:
$Rxy$ असणारा एक $y$ असेल, तर $f(x) = y$; अन्यथा $f(x) \fundefined$.
शिवाय $R$ हा \emph{प्रत्येक आदानाशी किमान एक प्रदान जोडणारा} असेल,
म्हणजे प्रत्येक $x \in A$ साठी $Rxy$ असणारा एक $y \in B$ असेल,
तर $f$ हे सर्वत्र परिभाषित असते.
\end{prop}

\begin{proof}
$Rxy$ असणारा एक $y$ आहे असे समजा. $Rxy'$ असणारा दुसरा
$y' \neq y$ असता, तर $R$ वरील अट मोडली गेली असती. म्हणून
$Rxy$ असणारा एक $y$ असेल, तर तो $y$ एकमेव असतो. त्यामुळे $f$
सुस्पष्टपणे परिभाषित आहे. उघडच $R_f = R$, आणि~$R$ प्रत्येक आदानाशी
किमान एक प्रदान जोडत असेल, तर $f$ हे सर्वत्र परिभाषित असते.
\end{proof}



